\chapter{Experimental Evaluation}
\label{ch:eval}

This chapter will focus on the experiments with the implementation of both the collision analysis and the subsumption analysis as it was described in Chapter~\ref{ch:impl}.
Additionally, this chapter will be closed with reasoning about the proposals for future work on this topic.
As the main contribution to the tool by this work is the encoding of \rego policies into SMT, which correspond to the step presented in Section~\ref{s:ir}, these experiments serve mainly as a demonstration of the tool's capabilities.

As each analysis targets a single tested variable, we will denote this variable as $x$ throughout the demonstration.
Every experiment will be marked by the input \rego policy, the output of the program with cosmetic edits, for the purpose of clear presentation, and potentially a quick note, regarding the experiment.

The symbolic context for these experiments is given by the request document:
\begin{center}
\texttt{input = \{ "count": $\Int$, "word": $\Str$ \}}
\end{center}

\section{Collision Analysis}

\begin{lstlisting}
package e01

x := 1
x := 2
\end{lstlisting}
\textit{collision on values \texttt{1} and \texttt{2}} \\
\texttt{input = \{ "count": 3, "word": "a" \}}

\noindent\rule{\linewidth}{0.4pt}
\begin{lstlisting}
package e02

x := 1 if input.count >= 1
x := 2 if input.count <= 1
\textit{collision on values \texttt{1} and \texttt{2}} \\
\texttt{input = \{ "count": 1, "word": "a" \}}

\noindent\rule{\linewidth}{0.4pt}
\begin{lstlisting}
package e03

x := true if input.word == "Hello, World!"
x := false if contains(input.word, "!")
\end{lstlisting}
\textit{collision on values \true and \false} \\
\texttt{input = \{ "count": 2, "word": "Hello, World!" \}}

\noindent\rule{\linewidth}{0.4pt}
\begin{lstlisting}
package e04

y := 1
x := z if z := y-1
x := y if input.count > 1
\end{lstlisting}
\textit{collision on values \texttt{z} and \texttt{y}} \\
\texttt{input = \{ "count": 2, "word": "c" \}}

\noindent\rule{\linewidth}{0.4pt}
\begin{lstlisting}
package e05

x := 1 if {
  input.count > 0
} else := 2 if {
  input.count < 10
} else := 3 if {
  input.count != 3
}

x := 3 if {
   input.count == 0
} else := 1 if {
  input.count == 1
}
\end{lstlisting}
\textit{collision on values \texttt{2} and \texttt{3}} \\
\texttt{input = \{ "count": 0, "word": "c" \}}

\noindent\rule{\linewidth}{0.4pt}
\begin{lstlisting}
package e06

x := 1 if input.count == 1
x := 2 if input.count == 2
x := input.count
\end{lstlisting}
\textit{no collisions}

In this example, $x$ is always assigned the value \texttt{input.count}; only the first two rule perform this assignment conditionally.

\noindent\rule{\linewidth}{0.4pt}
\begin{lstlisting}
package e07

x := 5
x := (2*y+10)/2-y

y := input.count
\end{lstlisting}
\textit{no collisions}

This example illustrates the encoding of the commonly known \emph{number forcing trick}, where any number maps to the same value.

%%%%%%%%%%%%%%%%%%%%%%%%%%%%%%%%%%%%%%%%%%%%
\section{Subsumption Analysis}

\begin{lstlisting}
package e01

x if input.count > 1
x if input.count > 10
\end{lstlisting}
\textit{subsumed condition: \texttt{input.count > 10}}

\noindent\rule{\linewidth}{0.4pt}
\begin{lstlisting}
package e02

x if {
  input.count > 1
} else if {
  input.count > 10
}
\end{lstlisting}
\textit{subsumed condition: \texttt{input.count > 10}}

In the case that the first condition fails, the second condition also always fails.

\noindent\rule{\linewidth}{0.4pt}
\begin{lstlisting}
package e03

x if contains(input.word, "X")
x if contains(input.word, "XX")
\end{lstlisting}
\textit{subsumed condition: \texttt{contains(input.word, "XX")}}

\noindent\rule{\linewidth}{0.4pt}
\begin{lstlisting}
package e04

y := input.count
x if {
  y > 5; y < 10
} else if {
  y > 0; y < 6
} else if {
  y > 3; y < 8
}
\end{lstlisting}
\textit{subsumed condition: \texttt{y > 3; y < 8}}

In this example, the subsumed condition is not directly implied by any of the other conditions, but rather with their disjunction.

\noindent\rule{\linewidth}{0.4pt}
\begin{lstlisting}
package e05

arr := [0,0,1,1,0,1]

x if arr[input.count] == 1
x if input.count == 3
x if input.count == 4
\end{lstlisting}
\textit{subsumed condition: \texttt{input.count == 3}}

In this example, the first condition asserts $x = \true$ when $\texttt{input.count} \in \{2,3,5\}$.

\section{Future Work}

Since \verirego is still in early development, there are multiple ways in which to expand this tool.
Firstly, \verirego does not support all \rego policies; we can enlarge the fragment of analysable policies by:
\begin{itemize}
  \item creating special SMT functions to represent more builtin functions,
  \item supporting more obscure \rego constructions, such as universal quantifiers,
  \item encoding accumulating rules.
\end{itemize}
Secondly, we can introduce more complex analysis, such as the comparison of multiple policies, i.e., \textit{one policy accepting a request implies another policy also accepting it}, or symbolic testing of a policy's properties.
Furthermore, the SMT encodings can be further optimized: either by choosing more fitted constructs for specific cases or by integrating more dedicated theory solvers.
